% =====================================================================
% Übungskapitel zu Kapitel 1 (Analysis) – steht direkt hinter Analysis.tex
% Nummerierung: Ü 1.n ; Lösungen im Lösungsteil am Ende des Buches (Loesungen_*.tex).
% =====================================================================
\Uebungskapitel{Übungen zur Analysis}

\Uebungsabschnitt{Lokales und globales Differenzieren}

\begin{uebung}*{Änderungsrate und h-Methode}
Gegeben ist die Funktion $f$ mit $f(x) = x^2-2x+3$ und $\mathds{D}_f = \mathds{R}$.
\begin{enumerate}[label=\alph*)]
  \item Berechnen Sie die mittlere Änderungsrate von $f$ im Intervall $[1;4]$.
  \item Bestimmen Sie $f'(2)$ mithilfe der h-Methode.
  \item Bestimmen Sie die Stelle $x_0$, an der die Tangente an $G_f$ parallel zur Sekante aus Teilaufgabe a) verläuft.
\end{enumerate}
\end{uebung}

\begin{uebung}*{Stetigkeit und Differenzierbarkeit}
\begin{enumerate}[label=\alph*)]
  \item Gegeben ist die Funktion $f$ mit $f(x) = |x-3|+1$ und $\mathds{D}_f = \mathds{R}$. Untersuchen Sie mithilfe der h-Methode, ob $f$ an der Stelle $x_0 = 3$ differenzierbar ist, und beschreiben Sie die Bedeutung Ihres Ergebnisses für den Graphen $G_f$.
  \item Begründen Sie, dass $f$ an der Stelle $x_0 = 3$ stetig ist.
  \item Gegeben ist die Funktion $g$ mit
  $$g(x)= \begin{cases}
              x^2 & \text{für}\hspace{0.5cm} x\leq 2\\
              a\cdot x+b & \text{für}\hspace{0.5cm}  x> 2
         \end{cases}\qquad a,b\in\mathds{R}.$$
  Bestimmen Sie $a$ und $b$ so, dass $g$ an der Stelle $x_0=2$ stetig und differenzierbar ist.
\end{enumerate}
\end{uebung}

\begin{uebung}*{Lineare Transformation}
\begin{enumerate}[label=\alph*)]
  \item Beschreiben Sie, wie der Graph der Funktion $g$ mit $g(x) = -2\cdot(x+1)^2+3$ aus der Normalparabel $y=x^2$ hervorgeht.
  \item Der Graph der Funktion $f$ mit $f(x) = \sqrt{x}$ wird zuerst an der $y$-Achse gespiegelt und anschließend um $2$ in Richtung der $y$-Achse nach oben verschoben. Geben Sie einen Funktionsterm des entstehenden Graphen sowie dessen Definitionsbereich an.
  \item Der Graph der Funktion $h$ mit $h(x) = \sin{(x)}$ wird in Richtung der $x$-Achse mit dem Faktor $\dfrac{1}{2}$ gestaucht. Geben Sie einen zugehörigen Funktionsterm sowie die Periode des neuen Graphen an.
\end{enumerate}
\end{uebung}

\begin{uebung}{Newton-Verfahren}
Gegeben ist die Funktion $f$ mit $f(x) = x^3+x-1$ und $\mathds{D}_f = \mathds{R}$.
\begin{enumerate}[label=\alph*)]
  \item Begründen Sie, dass $f$ genau eine Nullstelle besitzt und dass diese im Intervall $[0;1]$ liegt.
  \item Führen Sie ausgehend vom Startwert $x_0 = 1$ zwei Schritte des Newton-Verfahrens durch. Geben Sie $x_1$ und $x_2$ auf vier Nachkommastellen gerundet an.
  \item Für die Funktion $g$ mit $g(x)= x^3-3x$ soll eine Nullstelle mit dem Newton-Verfahren bestimmt werden. Erläutern Sie, warum der Startwert $x_0 = 1$ dafür ungeeignet ist.
\end{enumerate}
\end{uebung}

\Uebungsabschnitt{Kurvendiskussion}

\begin{uebung}*{Tangente und Normale}
Gegeben ist die Funktion $f$ mit $f(x) = x^3-3x^2+4$ und $\mathds{D}_f = \mathds{R}$.
\begin{enumerate}[label=\alph*)]
  \item Bestimmen Sie die Gleichung der Tangente an $G_f$ im Punkt $P(1|f(1))$.
  \item Bestimmen Sie die Gleichung der Normalen an $G_f$ im Punkt $P$.
  \item Ermitteln Sie die Koordinaten aller Punkte von $G_f$, in denen die Tangente parallel zur Geraden mit der Gleichung $y = 9x$ verläuft.
\end{enumerate}
\end{uebung}

\begin{uebung}*{Monotonie, Extrema und Krümmung}
Gegeben ist die Funktion $f$ mit $f(x) = x^3-6x^2+9x$ und $\mathds{D}_f = \mathds{R}$.
\begin{enumerate}[label=\alph*)]
  \item Bestimmen Sie die Nullstellen von $f$ und geben Sie deren Vielfachheit an.
  \item Untersuchen Sie das Monotonieverhalten von $G_f$ mithilfe einer Monotonietabelle und bestimmen Sie Art und Koordinaten der Extrempunkte.
  \item Bestimmen Sie das Krümmungsverhalten von $G_f$ sowie die Koordinaten des Wendepunkts.
  \item Bestimmen Sie die Gleichung der Wendetangente.
  \item Skizzieren Sie $G_f$ unter Verwendung der bisherigen Ergebnisse im Bereich $-0{,}5\leq x\leq 4{,}5$.
\end{enumerate}
\end{uebung}

\begin{uebung}*{Ausführliche Kurvendiskussion}
Gegeben ist die Funktion $f$ mit $f(x) = \dfrac{1}{4}x^4-x^3$ und $\mathds{D}_f = \mathds{D}_{\text{max}}$.
\begin{enumerate}[label=\alph*)]
  \item Geben Sie $\mathds{D}_{\text{max}}$ an und untersuchen Sie $G_f$ auf Symmetrie zum Koordinatensystem.
  \item Untersuchen Sie das Verhalten von $f$ an den Rändern des Definitionsbereichs und bestimmen Sie die gemeinsamen Punkte von $G_f$ mit den Koordinatenachsen.
  \item Bestimmen Sie Art und Lage aller Punkte von $G_f$ mit waagrechter Tangente. Begründen Sie insbesondere, dass im Ursprung ein Terrassenpunkt vorliegt.
  \item Bestimmen Sie die Wendepunkte von $G_f$.
  \item Geben Sie die Wertemenge von $f$ an und skizzieren Sie $G_f$.
\end{enumerate}
\end{uebung}

\begin{uebung}*{Graph der Ableitungsfunktion}
Die Abbildung zeigt den Graphen $G_{f'}$ der Ableitungsfunktion $f'$ einer in $\mathds{R}$ definierten ganzrationalen Funktion $f$. Der Graph $G_{f'}$ ist eine Parabel mit den Nullstellen $x=-1$ und $x=3$.
\begin{center}
\begin{tikzpicture}
\begin{axis}[ticklabel style={fill=none}, every tick label/.append style={fill=none},
        xmin=-2.5, xmax=4.5, ymin=-2.5, ymax=4,
        axis lines = middle,
        x=1cm, y=1cm,
        xtick={-2,-1,...,4},
        ytick={-2,-1,...,3},
        xlabel = $x$,
        ylabel = $y$]
  \addplot[color=black, thick, samples=200, domain=-2.2:4.2]{0.5*(x+1)*(x-3)};
  \draw[color=black] (4,2.6) node[left] {$G_{f'}$};
\end{axis}
\end{tikzpicture}
\end{center}
\begin{enumerate}[label=\alph*)]
  \item Geben Sie die Extremstellen von $f$ an und entscheiden Sie jeweils über die Art des Extremums. Begründen Sie Ihre Antwort.
  \item Geben Sie die Wendestelle von $f$ an und begründen Sie Ihre Angabe mithilfe des Graphen.
  \item Entscheiden Sie, ob $f(3) < f(0)$ gilt, und begründen Sie Ihre Entscheidung.
  \item Skizzieren Sie einen möglichen Graphen von $f$ mit $f(0)=0$.
\end{enumerate}
\end{uebung}

\Uebungsabschnitt{Stammfunktion}

\begin{uebung}*{Stammfunktionen ganzrationaler Funktionen}
\begin{enumerate}[label=\alph*)]
  \item Geben Sie jeweils eine Stammfunktion an:
  \[f(x) = 3x^2-4x+1,\qquad g(x) = \dfrac{1}{2}x^5 - x,\qquad h(x) = a^2x^3+2a \text{ mit } a\in\mathds{R}\]
  \item Bestimmen Sie diejenige Stammfunktion $F$ von $f$ mit $f(x) = 6x^2-2$, deren Graph durch den Punkt $P(1|3)$ verläuft.
  \item Begründen Sie, dass mit einer Stammfunktion $F$ von $f$ auch jede Funktion $G$ mit $G(x) = F(x)+c$ und $c\in\mathds{R}$ eine Stammfunktion von $f$ ist, und beschreiben Sie, wie die Graphen von $F$ und $G$ zueinander liegen.
\end{enumerate}
\end{uebung}

\Uebungsabschnitt{Funktionenscharen}

\begin{uebung}*{Funktionenschar und Ortskurve}
Gegeben ist die Funktionenschar $f_k$ mit $f_k(x) = x^3-k\cdot x$, $\mathds{D}_{f_k} = \mathds{R}$ und $k\in\mathds{R}$.
\begin{enumerate}[label=\alph*)]
  \item Untersuchen Sie in Abhängigkeit von $k$ die Anzahl der Nullstellen von $f_k$.
  \item Zeigen Sie, dass jeder Graph $G_{f_k}$ den Ursprung als Wendepunkt besitzt, und geben Sie die Steigung der Wendetangente in Abhängigkeit von $k$ an.
  \item Bestimmen Sie für $k>0$ die Extremstellen von $f_k$ und entscheiden Sie über die Art der Extrema.
  \item Bestimmen Sie eine Gleichung der Ortskurve, auf der die Hochpunkte aller Graphen $G_{f_k}$ mit $k>0$ liegen.
\end{enumerate}
\end{uebung}

\Uebungsabschnitt{Untersuchung von speziellen Funktionen}

\begin{uebung}*{Gebrochen-rationale Funktion}
Gegeben ist die Funktion $f$ mit $f(x) = \dfrac{x^2-4x+3}{x-2}$ und $\mathds{D}_f = \mathds{D}_{\text{max}}$.
\begin{enumerate}[label=\alph*)]
  \item Geben Sie $\mathds{D}_{\text{max}}$ an und bestimmen Sie die Nullstellen von $f$.
  \item Zeigen Sie, dass $f(x) = x-2-\dfrac{1}{x-2}$ gilt.
  \item Untersuchen Sie das Verhalten von $f$ an der Definitionslücke sowie für $x\longrightarrow\pm\infty$ und geben Sie die Gleichungen aller Asymptoten von $G_f$ an.
  \item Zeigen Sie, dass $G_f$ in jedem Teilintervall des Definitionsbereichs streng monoton wachsend ist.
  \item Begründen Sie, dass $G_f$ keinen Wendepunkt besitzt, obwohl sich das Krümmungsverhalten an der Stelle $x=2$ ändert.
  \item Skizzieren Sie $G_f$ zusammen mit seinen Asymptoten.
\end{enumerate}
\end{uebung}

\begin{uebung}*{Definitionslücken}
Gegeben ist die Funktion $g$ mit $g(x) = \dfrac{x^2-1}{x^2+x}$ und $\mathds{D}_g = \mathds{D}_{\text{max}}$.
\begin{enumerate}[label=\alph*)]
  \item Bestimmen Sie die Definitionslücken von $g$ und untersuchen Sie, ob es sich jeweils um eine Polstelle (mit oder ohne Vorzeichenwechsel) handelt oder ob die Definitionslücke durch Kürzen behoben werden kann.
  \item Geben Sie die Gleichung der waagrechten Asymptote von $G_g$ an und begründen Sie Ihre Angabe.
\end{enumerate}
\end{uebung}

\begin{uebung}*{Umkehrfunktion und Wurzelfunktionen}
\begin{enumerate}[label=\alph*)]
  \item Gegeben ist die Funktion $f$ mit $f(x) = (x-1)^2+2$ und $\mathds{D}_f = [1;\infty[$. Begründen Sie, dass $f$ umkehrbar ist, und bestimmen Sie einen Term der Umkehrfunktion $f^{-1}$ sowie deren Definitions- und Wertemenge.
  \item Der Graph der Funktion $g$ mit $g(x) = \sqrt{9-(x-2)^2}$ ist ein Halbkreis. Geben Sie den Mittelpunkt und den Radius des Halbkreises sowie $\mathds{D}_{\text{max}}$ von $g$ an.
  \item Gegeben ist die Funktion $h$ mit $h(x) = \sqrt{4x+1}$ und $\mathds{D}_h = \left[-\dfrac{1}{4};\infty\right[$. Bestimmen Sie die Gleichung der Tangente an $G_h$ im Punkt $Q(2|h(2))$.
\end{enumerate}
\end{uebung}

\begin{uebung}*{Trigonometrische Funktion}
Gegeben ist die Funktion $f$ mit $f(x) = 2\cos{(2x)}+1$ und $\mathds{D}_f = [0;\pi]$.
\begin{enumerate}[label=\alph*)]
  \item Geben Sie die Amplitude, die Periode und die Wertemenge von $f$ an.
  \item Bestimmen Sie die Nullstellen von $f$.
  \item Bestimmen Sie Art und Lage der Extrempunkte von $G_f$ (einschließlich der Randpunkte).
  \item Bestimmen Sie die Gleichung der Tangente an $G_f$ an der Stelle $x=\dfrac{\pi}{4}$.
\end{enumerate}
\end{uebung}

\begin{uebung}*{Exponentialfunktion}
Gegeben ist die Funktion $f$ mit $f(x) = (x+1)\cdot e^{-x}$ und $\mathds{D}_f = \mathds{R}$.
\begin{enumerate}[label=\alph*)]
  \item Bestimmen Sie die Nullstelle von $f$ sowie den Schnittpunkt von $G_f$ mit der $y$-Achse.
  \item Untersuchen Sie das Verhalten von $f$ für $x\longrightarrow\pm\infty$ und geben Sie die Gleichung der waagrechten Asymptote an.
  \item Zeigen Sie, dass $f'(x) = -x\cdot e^{-x}$ gilt, und bestimmen Sie Art und Lage des Extrempunkts von $G_f$.
  \item Bestimmen Sie die Koordinaten des Wendepunkts von $G_f$.
  \item Zeigen Sie, dass $F$ mit $F(x) = -(x+2)\cdot e^{-x}$ eine Stammfunktion von $f$ ist.
\end{enumerate}
\end{uebung}

\begin{uebung}*{Logarithmusfunktion}
Gegeben ist die Funktion $f$ mit $f(x) = \ln{(x^2+1)}$ und $\mathds{D}_f = \mathds{R}$.
\begin{enumerate}[label=\alph*)]
  \item Begründen Sie, dass $\mathds{D}_f = \mathds{R}$ der maximale Definitionsbereich ist und dass $G_f$ achsensymmetrisch zur $y$-Achse ist.
  \item Bestimmen Sie Art und Lage des Extrempunkts von $G_f$.
  \item Zeigen Sie, dass $f''(x) = \dfrac{2\cdot(1-x^2)}{(x^2+1)^2}$ gilt, und bestimmen Sie die Wendepunkte von $G_f$.
  \item Bestimmen Sie die Gleichung der Wendetangente im Wendepunkt mit positiver $x$-Koordinate.
\end{enumerate}
\end{uebung}

\Uebungsabschnitt{Anwendungen der Differentialrechnung}

\begin{uebung}*{Modellierung von Funktionen}
\begin{enumerate}[label=\alph*)]
  \item Der Graph einer ganzrationalen Funktion dritten Grades ist punktsymmetrisch zum Ursprung und besitzt den Tiefpunkt $T(1|-2)$. Bestimmen Sie einen Funktionsterm.
  \item Eine Parabel hat den Scheitelpunkt $S(2|-2)$ und schneidet die $y$-Achse im Punkt $(0|2)$. Bestimmen Sie den zugehörigen Funktionsterm in Scheitelpunktform und in allgemeiner Form.
\end{enumerate}
\end{uebung}

\begin{uebung}*{Extremwertproblem}
Gegeben ist die Funktion $f$ mit $f(x) = 9-x^2$ und $\mathds{D}_f=\mathds{R}$. Einem Parabelsegment, das von $G_f$ und der $x$-Achse eingeschlossen wird, soll ein Rechteck einbeschrieben werden. Zwei Eckpunkte des Rechtecks liegen auf der $x$-Achse, die beiden anderen Eckpunkte $P(u|f(u))$ und $Q(-u|f(u))$ mit $0<u<3$ auf $G_f$.
\begin{enumerate}[label=\alph*)]
  \item Zeigen Sie, dass sich der Flächeninhalt des Rechtecks durch $A(u) = 18u-2u^3$ beschreiben lässt.
  \item Bestimmen Sie $u$ so, dass der Flächeninhalt maximal wird, und geben Sie den maximalen Flächeninhalt an.
  \item Begründen Sie, dass an den Rändern des Definitionsbereichs von $A$ kein größerer Flächeninhalt auftreten kann.
\end{enumerate}
\end{uebung}

\Uebungsabschnitt{Bestimmtes Integral, Integralfunktion und Integrationsregeln}

\begin{uebung}*{Bestimmte Integrale}
Berechnen Sie die folgenden Integrale.
\begin{multicols}{2}
\begin{enumerate}[label=\alph*)]
  \item $\integral[limits-mode=limits,limits={0, 2}]{(3x^2-2x)}$
  \item $\integral[limits-mode=limits,limits={1, 4}]{\dfrac{1}{\sqrt{x}}}$
  \item $\integral[limits-mode=limits,limits={0, \pi}]{\sin{(x)}}$
  \item $\integral[limits-mode=limits,limits={0, \ln{(2)}}]{e^{2x}}$
  \item $\integral[limits-mode=limits,limits={1, e}]{\dfrac{1}{x}}$
  \item $\integral[limits-mode=limits,limits={0, 1}]{2x\cdot e^{x^2}}$
  \item $\integral[limits-mode=limits,limits={0, 1}]{\dfrac{2x}{x^2+1}}$
  \item $\integral[limits-mode=limits,limits={1, 2}]{\ln{(x)}}$
\end{enumerate}
\end{multicols}
\end{uebung}

\begin{uebung}*{Integralfunktion und Flächenbilanz}
Gegeben ist die Integralfunktion $I_1$ mit $I_1(x) = \integral[limits-mode=limits,limits={1, x}, variable={t}]{(2t-4)}$ und $x\in\mathds{R}$.
\begin{enumerate}[label=\alph*)]
  \item Bestimmen Sie einen integralfreien Term von $I_1$ sowie die Nullstellen von $I_1$.
  \item Begründen Sie ohne Rechnung mithilfe des Hauptsatzes der Differential- und Integralrechnung, dass $I_1$ an der Stelle $x=2$ einen Tiefpunkt besitzt.
  \item Berechnen Sie $\integral[limits-mode=limits,limits={0, \frac{\pi}{2}}]{\cos{(2x)}}$ und deuten Sie das Ergebnis als Flächenbilanz.
\end{enumerate}
\end{uebung}

\Uebungsabschnitt{Flächenberechnungen, uneigentliches Integral und Rotationskörper}

\begin{uebung}*{Flächeninhalte}
\begin{enumerate}[label=\alph*)]
  \item Gegeben ist die Funktion $f$ mit $f(x) = x^3-4x$ und $\mathds{D}_f = \mathds{R}$. Berechnen Sie den Inhalt der Fläche, die $G_f$ mit der $x$-Achse einschließt, und erläutern Sie, warum $\integral[limits-mode=limits,limits={-2, 2}]{f(x)}$ nicht diesen Flächeninhalt angibt.
  \item Gegeben sind die Funktionen $p$ mit $p(x) = x^2$ und $q$ mit $q(x) = 2x+3$, jeweils mit $\mathds{D} = \mathds{R}$. Berechnen Sie den Inhalt der Fläche, die von den Graphen $G_p$ und $G_q$ eingeschlossen wird.
\end{enumerate}
\end{uebung}

\begin{uebung}*{Uneigentliches Integral und Rotationskörper}
Gegeben ist die Funktion $f$ mit $f(x) = \dfrac{1}{x}$ und $\mathds{D}_f = [1;\infty[$.
\begin{enumerate}[label=\alph*)]
  \item Zeigen Sie, dass für den Inhalt $A(z)$ der Fläche zwischen $G_f$ und der $x$-Achse über dem Intervall $[1;z]$ mit $z>1$ gilt: $A(z) = \ln{(z)}$. Untersuchen Sie, ob die ins Unendliche reichende Fläche einen endlichen Inhalt besitzt.
  \item Die Fläche zwischen $G_f$ und der $x$-Achse über dem Intervall $[1;z]$ rotiert um die $x$-Achse. Bestimmen Sie das Volumen $V(z)$ des entstehenden Rotationskörpers sowie $\lim\limits_{z\longrightarrow\infty} V(z)$.
  \item Beurteilen Sie die Ergebnisse aus a) und b) im Vergleich.
  \item Die Fläche zwischen dem Graphen der Funktion $w$ mit $w(x) = \sqrt{x}$ und der $x$-Achse über dem Intervall $[0;4]$ rotiert um die $x$-Achse. Berechnen Sie das Volumen des entstehenden Rotationskörpers.
\end{enumerate}
\end{uebung}

\Uebungsabschnitt{Das Integral als Gesamtänderung einer Größe}

\begin{uebung}{Zufluss in ein Becken}
In ein Becken, das zu Beginn der Beobachtung $10\ \text{m}^3$ Wasser enthält, fließt Wasser. Die Zuflussrate wird für $0\leq t\leq 8$ durch die Funktion $r$ mit $r(t) = -\dfrac{1}{2}t^2+4t$ beschrieben ($t$ in Stunden seit Beobachtungsbeginn, $r(t)$ in $\frac{\text{m}^3}{\text{h}}$). Es fließt kein Wasser ab.
\begin{enumerate}[label=\alph*)]
  \item Bestimmen Sie den Zeitpunkt, zu dem die Zuflussrate maximal ist, und geben Sie die maximale Zuflussrate an.
  \item Interpretieren Sie den Term $\integral[limits-mode=limits,limits={2, 4}, variable={t}]{r(t)}$ im Sachzusammenhang und berechnen Sie seinen Wert.
  \item Bestimmen Sie einen Term für die Funktion $V$, die das Wasservolumen im Becken in Abhängigkeit von der Zeit $t$ angibt, und berechnen Sie das Wasservolumen nach $8$ Stunden.
  \item Begründen Sie, dass das Wasservolumen im Becken im gesamten Beobachtungszeitraum zunimmt, die Zunahme aber ab $t=4$ immer langsamer erfolgt.
\end{enumerate}
\end{uebung}

\Uebungsabschnitt{Abiturniveau \textendash\ Teil A (ohne Hilfsmittel)}

\begin{uebung}*{Funktion mit Logarithmus}
Gegeben ist die Funktion $f$ mit $f(x) = x\cdot\ln{(x)}$ und $\mathds{D}_f = \mathds{R}^+$.
\begin{enumerate}[label=\alph*)]
  \item Bestimmen Sie die Nullstelle von $f$.
  \item Zeigen Sie, dass $f'(x) = \ln{(x)}+1$ gilt, und bestimmen Sie Art und Lage des Extrempunkts von $G_f$.
  \item Zeigen Sie, dass $F$ mit $F(x) = \dfrac{1}{2}x^2\cdot\ln{(x)}-\dfrac{1}{4}x^2$ eine Stammfunktion von $f$ ist.
\end{enumerate}
\end{uebung}

\begin{uebung}*{Funktionsterme mit vorgegebenen Eigenschaften}
Geben Sie jeweils den Term einer Funktion an, die die angegebenen Eigenschaften besitzt.
\begin{enumerate}[label=\alph*)]
  \item Der Graph besitzt die senkrechte Asymptote $x=2$ und die waagrechte Asymptote $y=1$.
  \item Die Funktion hat den maximalen Definitionsbereich $]-\infty;3]$.
  \item Die Funktion ist periodisch mit der Periode $4$ und hat die Wertemenge $[1;3]$.
\end{enumerate}
\end{uebung}

\begin{uebung}*{Integrale}
\begin{enumerate}[label=\alph*)]
  \item Berechnen Sie $\integral[limits-mode=limits,limits={0, 2}]{e^{\frac{1}{2}x}}$.
  \item Bestimmen Sie $b>0$ so, dass $\integral[limits-mode=limits,limits={0, b}]{(2x-2)} = 0$ gilt.
  \item Deuten Sie das Ergebnis aus b) geometrisch.
\end{enumerate}
\end{uebung}

\begin{uebung}*{Exponentialfunktion mit Substitution}
Gegeben ist die Funktion $g$ mit $g(x) = e^{2x}-4e^x+3$ und $\mathds{D}_g = \mathds{R}$.
\begin{enumerate}[label=\alph*)]
  \item Bestimmen Sie die Nullstellen von $g$. (Hinweis: Setzen Sie $z = e^x$.)
  \item Bestimmen Sie Art und Lage des Extrempunkts von $G_g$.
  \item Untersuchen Sie das Verhalten von $g$ für $x\longrightarrow-\infty$ und geben Sie die Gleichung der waagrechten Asymptote an.
\end{enumerate}
\end{uebung}

\begin{uebung}*{Tangenten von einem äußeren Punkt}
Gegeben ist die Funktion $f$ mit $f(x) = x^2$ und $\mathds{D}_f = \mathds{R}$ sowie der Punkt $P(1|-3)$, der nicht auf $G_f$ liegt.
\begin{enumerate}[label=\alph*)]
  \item Zeigen Sie, dass die Tangente an $G_f$ im Punkt $(a|a^2)$ die Gleichung $y = 2a\cdot x-a^2$ hat.
  \item Bestimmen Sie die Gleichungen aller Tangenten an $G_f$, die durch den Punkt $P$ verlaufen.
\end{enumerate}
\end{uebung}

\begin{uebung}*{Fläche zwischen Graph und Umkehrgraph}
Gegeben sind die Funktionen $w$ mit $w(x) = \sqrt{x}$ und $\mathds{D}_w = \mathds{R}_0^+$ sowie $\ell$ mit $\ell(x) = \dfrac{1}{2}x$ und $\mathds{D}_\ell = \mathds{R}$.
\begin{enumerate}[label=\alph*)]
  \item Bestimmen Sie die Schnittstellen von $G_w$ und $G_\ell$.
  \item Berechnen Sie den Inhalt der Fläche, die $G_w$ und $G_\ell$ einschließen.
  \item Spiegelt man beide Graphen an der Winkelhalbierenden $y=x$, entstehen die Graphen von $x\mapsto x^2$ ($x\geq 0$) und $x\mapsto 2x$. Begründen Sie ohne Rechnung, welchen Inhalt die von diesen Graphen eingeschlossene Fläche hat.
\end{enumerate}
\end{uebung}

\Uebungsabschnitt{Abiturniveau \textendash\ Teil B (mit Hilfsmitteln)}

\begin{uebung}{Medikament im Blut}
Die Konzentration eines Medikaments im Blut eines Patienten wird für $t\geq 0$ durch die Funktion $c$ mit $c(t) = 10t\cdot e^{-0{,}5t}$ beschrieben ($t$ in Stunden nach der Einnahme, $c(t)$ in $\frac{\text{mg}}{\text{l}}$).
\begin{enumerate}[label=\alph*)]
  \item Bestimmen Sie den Zeitpunkt, zu dem die Konzentration maximal ist, sowie die maximale Konzentration.
  \item Bestimmen Sie den Zeitpunkt, zu dem die Konzentration am stärksten abnimmt, und geben Sie die zugehörige Änderungsrate an.
  \item Das Medikament wirkt, solange die Konzentration mindestens $4\,\frac{\text{mg}}{\text{l}}$ beträgt. Bestimmen Sie die Dauer der Wirkung.
  \item Zeigen Sie, dass $C$ mit $C(t) = -(20t+40)\cdot e^{-0{,}5t}$ eine Stammfunktion von $c$ ist, und berechnen Sie die mittlere Konzentration in den ersten $6$ Stunden nach der Einnahme.
  \item Beschreiben Sie das langfristige Verhalten der Konzentration und begründen Sie Ihre Aussage.
\end{enumerate}
\end{uebung}

\begin{uebung}{Funktionenschar mit Exponentialfunktion}
Gegeben ist die Funktionenschar $f_a$ mit $f_a(x) = x\cdot e^{-ax}$, $a>0$ und $\mathds{D}_{f_a} = \mathds{R}$.
\begin{enumerate}[label=\alph*)]
  \item Bestimmen Sie die Nullstelle von $f_a$ und untersuchen Sie das Verhalten von $f_a$ für $x\longrightarrow\pm\infty$.
  \item Bestimmen Sie Art und Lage des Extrempunkts von $G_{f_a}$.
  \item Zeigen Sie, dass alle Hochpunkte auf der Geraden mit der Gleichung $y = \dfrac{1}{e}\cdot x$ liegen.
  \item Bestimmen Sie die Wendestelle von $f_a$.
  \item Zeigen Sie, dass $F_a$ mit $F_a(x) = -\left(\dfrac{x}{a}+\dfrac{1}{a^2}\right)\cdot e^{-ax}$ eine Stammfunktion von $f_a$ ist. Der Graph $G_{f_a}$ schließt für $x\geq 0$ mit der $x$-Achse eine ins Unendliche reichende Fläche ein. Zeigen Sie, dass diese den Inhalt $\dfrac{1}{a^2}$ hat, und bestimmen Sie $a$ so, dass der Inhalt $4$ beträgt.
\end{enumerate}
\end{uebung}

\begin{uebung}{Skateboard-Rampe}
Das Profil einer Skateboard-Rampe soll vom Boden im Punkt $O(0|0)$ zur oberen Plattform im Punkt $P(4|2)$ führen (Angaben in Metern). An beiden Punkten soll das Profil ohne Knick, also waagrecht, anschließen. Das Profil wird durch den Graphen einer ganzrationalen Funktion $f$ dritten Grades mit $0\leq x\leq 4$ modelliert.
\begin{enumerate}[label=\alph*)]
  \item Bestimmen Sie einen Term von $f$.\\
  {\small (zur Kontrolle: $f(x) = -\frac{1}{16}x^3+\frac{3}{8}x^2$)}
  \item Bestimmen Sie die Stelle, an der die Rampe am steilsten ist, und den zugehörigen Steigungswinkel.
  \item Die Rampe ist $3\,\text{m}$ breit und wird unterhalb des Profils vollständig mit Beton gefüllt. Berechnen Sie das benötigte Betonvolumen.
\end{enumerate}
\end{uebung}

\begin{uebung}{Tagestemperatur}
Der Temperaturverlauf an einem Sommertag wird für $0\leq t\leq 24$ durch die Funktion $T$ mit $T(t) = 18-6\cdot\cos{\left(\dfrac{\pi}{12}(t-3)\right)}$ beschrieben ($t$ in Stunden nach Mitternacht, $T(t)$ in $^{\circ}$C).
\begin{enumerate}[label=\alph*)]
  \item Bestimmen Sie die höchste und die niedrigste Temperatur sowie die zugehörigen Uhrzeiten.
  \item Ermitteln Sie, wie lange die Temperatur an diesem Tag mindestens $21\,^{\circ}$C beträgt.
  \item Bestimmen Sie den Zeitpunkt, zu dem die Temperatur am stärksten ansteigt, und die zugehörige Änderungsrate.
  \item Berechnen Sie die mittlere Temperatur des Tages und deuten Sie das Ergebnis am Funktionsterm.
\end{enumerate}
\end{uebung}

\begin{uebung}{Gebrochen-rationale Funktion und Fläche}
Gegeben ist die Funktion $f$ mit $f(x) = \dfrac{8x}{x^2+4}$ und $\mathds{D}_f = \mathds{R}$.
\begin{enumerate}[label=\alph*)]
  \item Begründen Sie, dass $\mathds{D}_f = \mathds{R}$ gilt, untersuchen Sie $G_f$ auf Symmetrie und geben Sie die Gleichung der Asymptote an.
  \item Bestimmen Sie Art und Lage der Extrempunkte von $G_f$.
  \item Es gilt $f''(x) = \dfrac{16x\cdot(x^2-12)}{(x^2+4)^3}$. Bestimmen Sie die Wendepunkte von $G_f$.
  \item $G_f$ schließt mit der $x$-Achse und der Geraden $x=u$ mit $u>0$ eine Fläche ein. Zeigen Sie, dass für ihren Inhalt $A(u) = 4\ln{\left(\dfrac{u^2+4}{4}\right)}$ gilt, und bestimmen Sie $u$ so, dass $A(u) = 4\ln{(5)}$ ist.
  \item Skizzieren Sie $G_f$ im Bereich $-6\leq x\leq 6$.
\end{enumerate}
\end{uebung}

\begin{uebung}{Sektglas}
Der Innenraum eines Sektglases entsteht, wenn der Graph der Funktion $r$ mit $r(x) = \sqrt{2x}$ für $0\leq x\leq 8$ um die $x$-Achse rotiert. Dabei ist $x$ die Höhe über dem tiefsten Punkt des Innenraums (in cm).
\begin{enumerate}[label=\alph*)]
  \item Zeigen Sie, dass das Volumen bei der Füllhöhe $h$ (in cm) $V(h) = \pi h^2$ (in $\text{cm}^3$) beträgt, und berechnen Sie das Volumen des vollen Glases.
  \item Bestimmen Sie die Füllhöhe, bei der das Glas genau zur Hälfte gefüllt ist, und vergleichen Sie mit der halben Glashöhe.
  \item Das Glas wird gleichmäßig mit $20\,\frac{\text{ml}}{\text{s}}$ gefüllt. Berechnen Sie die Zeit bis zum vollständigen Füllen und begründen Sie, dass der Flüssigkeitsspiegel zu Beginn schneller steigt als am Ende.
\end{enumerate}
\end{uebung}

\begin{uebung}{Logistisches Wachstum}
Die Anzahl der Bakterien in einer Nährlösung wird für $t\geq 0$ durch $B(t) = \dfrac{1000}{1+9e^{-0{,}5t}}$ beschrieben ($t$ in Stunden, $B(t)$ in Tausend).
\begin{enumerate}[label=\alph*)]
  \item Bestimmen Sie den Anfangsbestand und den langfristigen Bestand (Sättigungsgrenze).
  \item Berechnen Sie, nach welcher Zeit die Hälfte der Sättigungsgrenze erreicht ist.
  \item Zeigen Sie, dass $B'(t) = 0{,}5\cdot B(t)\cdot\left(1-\dfrac{B(t)}{1000}\right)$ gilt.
  \item Begründen Sie mithilfe von c), dass das Wachstum am stärksten ist, wenn der Bestand $500$ beträgt, und geben Sie die maximale Wachstumsrate an.
\end{enumerate}
\end{uebung}

\begin{uebung}{Schachtel mit maximalem Volumen}
Aus einem rechteckigen Karton mit den Maßen $30\,\text{cm}\times 16\,\text{cm}$ wird eine oben offene Schachtel hergestellt. Dazu werden an den vier Ecken gleich große Quadrate mit der Seitenlänge $x$ (in cm) ausgeschnitten und die Seiten hochgeklappt.
\begin{enumerate}[label=\alph*)]
  \item Stellen Sie das Volumen $V(x)$ der Schachtel auf und geben Sie einen sinnvollen Definitionsbereich an.
  \item Bestimmen Sie $x$ so, dass das Volumen maximal wird, und geben Sie das maximale Volumen an.
\end{enumerate}
\end{uebung}

\begin{uebung}{Tunnelquerschnitt}
Für den Querschnitt eines Tunnels werden zwei Modelle verglichen ($x$ und $y$ in m, Fahrbahn auf der $x$-Achse):
\[p(x) = 4-\dfrac{1}{4}x^2\qquad\text{und}\qquad k(x) = 4\cos{\left(\dfrac{\pi}{8}x\right)},\qquad -4\leq x\leq 4.\]
\begin{enumerate}[label=\alph*)]
  \item Zeigen Sie, dass beide Modelle dieselbe Breite am Boden und dieselbe maximale Höhe haben.
  \item Berechnen Sie für beide Modelle die Querschnittsfläche.
  \item Ein Fahrzeug ist $2{,}5\,\text{m}$ breit und $3\,\text{m}$ hoch und fährt genau in der Mitte. Prüfen Sie, ob es in beiden Modellen durch den Tunnel passt.
  \item Bestimmen Sie für beide Modelle die Breite des Tunnels in $3\,\text{m}$ Höhe.
\end{enumerate}
\end{uebung}

\UebungskapitelEnde
